\documentclass[12pt]{article}
\usepackage[utf8]{inputenc}
\usepackage[T1]{fontenc}
\usepackage{hyperref}
\title{Tumshuqese Language Resources for K-12 Curriculum Development}
\author{Compiled by Claude Sonnet 5.5 (claude-sonnet-5-5) for Andr\'as Kornai}
\date{}
\begin{document}
\maketitle

\section{Introduction}
Tumshuqese (ISO 639-3 xtq, Glottocode tums1237) is an extinct Eastern Iranian language, a dialect of Saka closely related to Khotanese, known from manuscripts from Tumshuq on the northern edge of the Tarim Basin, mostly of the late 7th or 8th century. It is much less understood than Khotanese because fewer texts survive. There are no speakers and no school programme; none was found in the searches.

\section{Primary Educational Resources}
None found.

\section{Linguistic and Reference Materials}
\begin{itemize}
\item Wikipedia, ``Saka language'' (covers Khotanese and Tumshuqese): \url{https://en.wikipedia.org/wiki/Saka_language}
\item Ogihara Hirotoshi, ``Miscellany on the Tumshuqese Documents'', \emph{Central Asiatic Journal} (2020), on contracts and fixed expressions: \url{https://caj.harrassowitz-library.com/article/CAJ/2020/1-2/3}
\item ``Miscellany on the Tumshuqese Documents: Part II'': \url{https://journals.eco-vector.com/1811-8062/article/view/77877}
\item Skj\ae rv\o, P.\,O., ``On the Tumshuqese Karmav\=ac\=an\=a Text'' (\emph{Journal of the Royal Asiatic Society}): \url{https://resolve.cambridge.org/core/journals/journal-of-the-royal-asiatic-society/article/on-the-tumshuqese-karmavacana-text1/A4D0D647309421D71555FEC3DA2CC092}
\end{itemize}
No dictionary or grammar of Tumshuqese was found in the search results.

\section{Religious and Community Resources}
The texts include a Buddhist ritual text (the Karmav\=ac\=an\=a) and secular contracts. No living community was found.

\section{Digital and Audio Resources}
None found. The language has no audio.

\section{Curriculum Development Context}
Instruction in Tumshuqese is not feasible. Possible enrichment themes are Silk Road history, the Brahmi script and Iranian languages of Central Asia, drawing on the scholarly articles above with specialist help.
\end{document}